% Generated draft: requires mathematical and source review.
\documentclass[11pt]{article}
\usepackage[margin=25mm]{geometry}
\usepackage{fontspec}
\setmainfont{Latin Modern Roman}
\usepackage{amsmath,amssymb,mathtools,unicode-math}
\setmathfont{Latin Modern Math}
\usepackage{enumitem,needspace,xurl,hyperref}
\setlength{\parindent}{0pt}
\setlength{\parskip}{5pt}
\setlength{\emergencystretch}{3em}
\newcommand{\N}{\mathbb{N}}
\newcommand{\Z}{\mathbb{Z}}
\newcommand{\Q}{\mathbb{Q}}
\newcommand{\R}{\mathbb{R}}
\newcommand{\C}{\mathbb{C}}
\newcommand{\F}{\mathbb{F}}
\newcommand{\A}{\mathbb{A}}
\newcommand{\PP}{\mathbb{P}}
\newcommand{\E}{\mathbb{E}}
\newcommand{\eps}{\varepsilon}
\newcommand{\dd}{\,\mathrm{d}}
\begin{document}
\section*{UnsolvedMath ID: 20002464}
\textbf{Source problem code:} catalogue entry
\subsection{Definitions and mathematical statement}
The source is Theodoros Katsaounis's research contribution on relaxation methods for balance laws. It discusses linear hyperbolic relaxation systems with stiff lower-order terms, avoiding Riemann solvers and characteristic decompositions at the cost of additional unknowns. Its numerical topics include first- and second-order finite-volume shallow-water schemes in one and two dimensions, two source-term treatments, dry beds, steady-state preservation to relaxation-parameter accuracy, passive pollutant transport/diffusion, and finite-element/Runge--Kutta methods with adaptive shock resolution and low-, high- and mixed-order $h$--$p$ variants. The printed relaxation-parameter symbol is damaged, and no exact error inequality is supplied.

The final paragraph asks a genuine additional question: in thermomechanical shear-band formation, do temperature and strain rate become unbounded in finite or only infinite time, and how can the singular regime be simulated accurately? The source describes a delta-like velocity, whereas the aid says the cited [BKT] account instead has delta-like strain rate $v_x$ and a step-like velocity $v$. These distinct statements must not be conflated. The complete constitutive law, initial and boundary data, parameter regime and solution class are absent from the original contribution.

The aid's energy diagnostic can be expressed for the explicitly normalized model on $0<x<1$:
\[
 v_t=\sigma_x,\qquad \gamma_t=v_x,
 \qquad\theta_t=\sigma v_x,\qquad
 \sigma=\theta^{-\alpha}\gamma^m(v_x)^n,
\]
with positive temperature $\theta$, positive strain $\gamma$, positive strain rate $v_x$, and specified exponents. This power-law system is an illustrative research-aid setting, not a recovered equation from the original paragraph. For a classical solution its exact boundary-work identity is
\[
 \frac{\dd}{\dd t}\int_0^1\left(\tfrac12v^2+\theta\right)\,\dd x
       =[v\sigma]_{x=0}^{x=1}.
\]
If the cumulative positive boundary work up to a finite time $T$ is bounded by $W_*$, then for $0\le t<T$,
\[
 \tfrac12\|v(t)\|_2^2+\|\theta(t)\|_1\le C,
 \qquad C=\int_0^1(\tfrac12v_0^2+\theta_0)\,\dd x+W_*.
\]
Consequently no nonzero Dirac mass can be a distributional velocity limit along that interval, because the velocity remains bounded in $L^2$. A step limit for $v$, with an atomic derivative $v_x$, is not excluded. Also
\[
 |\{x:\theta(t,x)\ge H\}|\le C/H\qquad(H>0),
\]
so any arbitrarily high temperature values must occupy shrinking superlevel sets. These bounds do not exclude finite-time $L^\infty$ concentration of temperature or strain rate.

For an exact benchmark choose $\alpha>0$, $m\ge0$, $n>0$, and positive constants $\gamma_0,\theta_0$. Uniform shear has
\[
 v=x,\qquad\gamma=t+\gamma_0,\qquad
 \theta(t)=\left[\theta_0^{\alpha+1}
 +\frac{\alpha+1}{m+1}
  \bigl((t+\gamma_0)^{m+1}-\gamma_0^{m+1}\bigr)\right]^{1/(\alpha+1)}.
\]
It exists for all finite times and its temperature grows only as $t\to\infty$. Its tangent momentum diffusivity is
\[
 \partial_{v_x}\sigma
   =n\theta(t)^{-\alpha}(t+\gamma_0)^m
   \asymp t^{(m-\alpha)/(\alpha+1)}
 \quad(t\to\infty),
\]
which collapses when thermal softening dominates hardening, $\alpha>m$. Thus diffusion collapse is compatible with infinite-time rather than finite-time growth in this exact family. The benchmark and finite-work constraints do not settle the original shear-band blow-up question for general data.
\subsection{Short English statement}
Preserve the relaxation-method agenda and its shear-band question about finite- versus infinite-time blow-up. The supplied aid distinguishes a strain-rate delta from a velocity delta, gives a finite-boundary-work concentration constraint and an exact uniform-shear benchmark, and leaves general blow-up unresolved.
\subsection{Sources}
\begin{itemize}
\item Theodoros Katsaounis, contribution A.11, American Institute of Mathematics workshop document \emph{Stiff Sources and Numerical Methods for Conservation Laws}, version 1 April 2005, \url{https://aimath.org/WWN/balancelaws/balancelaws.pdf}.
\item Source reference labels [JX], [DK1], [DK2], [DK3], [AKM], [KM], [GM], [AMT] and [BKT]. Their full bibliographic entries are not included in the supplied record; [BKT] is the shear-band reference whose strain-rate/velocity description is compared in the aid.
\item Local UnsolvedMath registry, permanent ID 20002464, source code \texttt{catalogue entry}, canonical location \texttt{aim-pdes-notes.json}, \texttt{notes[75]}, record dated 14 August 2026. Statement status is reconstructed unverified, clean statement null, and original literature field empty.
\item Supplied research attempt 1: partial result, problem status at run partially solved, rigor labelled proved, novelty confidence low. The boundary-work obstruction and uniform-shear diagnostic are machine-generated research aid, not peer reviewed; independently verify before citation.
\end{itemize}
\end{document}
